% !TeX root = euclid-without-words.tex
% !TeX Program=pdfLaTeX
%
% Perpendicular bisectors meet in the center of the circumscribed circle
%
\begin{minipage}[t]{.48\textwidth}
\begin{tikzpicture}[scale=.8,baseline=-30mm]
% Circle goes from (a) to (b)
\coordinate (a) at (0,0);
\coordinate (b) at (6,0);
% Line containing lower points is (c) -- (d)
\coordinate (c) at (0,-2);
\coordinate (d) at (6,-2);
% Line containing upper points is (e) -- (f)
\coordinate (e) at (0,2);
\coordinate (f) at (4,3);
% Name the two chords
\path [name path=chord1] (c) -- (d);
\path [name path=chord2] (e) -- (f);
% Name the coordinate of the center of the circle
\coordinate (center) at ($ (a)!.5!(b) $);
% Draw a circle whose center is half-way between (a) and (b) through (a)
\node [draw,thick,circle through=(a),name path=circ] at (center) {};
% Get intersections of the upper and lower lines with the circle
\path [name intersections={of=circ and chord1,by={i1,i2}}];
\path [name intersections={of=circ and chord2,by={i3,i4}}];
% Draw triangle
\draw [thick,blue] (i1) -- node[left] {$\mathbf{a}$} ($(i1)!.5!(i3)$) --
  node[left] {$\mathbf{a}$} (i3);
\draw[thick,black!10!teal] (i3) -- node[right] {$\mathbf{b}$} ($(i3)!.5!(i2)$) --
  node[right] {$\mathbf{b}$} (i2);
\draw[thick,red] (i2) -- node[below] {$\mathbf{c}$} ($(i2)!.5!(i1)$) -- 
  node[below] {$\mathbf{c}$} (i1);
% Draw three perpendicular bisectors
\draw [thick,red,fill] (center) -- ($(i1)!(center)!(i2)$);
\draw [thick,blue,fill] (center) -- ($(i1)!(center)!(i3)$);
\draw [thick,black!10!teal,fill] (center) -- ($(i2)!(center)!(i3)$);
% Draw squares to denote right angles
\draw[red] ($ (i1) !.5! (i2) $) rectangle +(8pt,8pt);
\draw[blue,rotate=-30] ($ (i1) !.5! (i3) $) rectangle +(8pt,8pt);
\draw[violet,rotate=-160] ($ (i3) !.5! (i2) $)
  rectangle +(8pt,8pt);
\end{tikzpicture}
\end{minipage}
